\chapter*{Résumé}
\addcontentsline{toc}{chapter}{Résumé}

\textbf{Mots-clés :} IoT, sécurité, TLS 1.3, mTLS, PKI, HashiCorp Vault, MQTT, Suricata, Wazuh, SIEM, IsolationForest, RandomForest, détection d'anomalies, GNS3, Scrum.

\vspace{0.5cm}

L'essor de l'Internet des Objets (IoT) dans les environnements des opérateurs de télécommunications engendre des flux de données massifs et hétérogènes exposés à des risques croissants d'interception, d'altération et de compromission. Ce projet de fin d'études, réalisé en partenariat avec \TT{}, aborde la problématique de la \textbf{sécurisation de bout en bout des flux de communication d'un écosystème IoT simulé}.

La démarche adoptée suit la méthodologie \textbf{Scrum} et s'appuie sur un laboratoire simulé combinant \textbf{GNS3} pour la topologie réseau et \textbf{Docker} pour les nœuds applicatifs. Nous avons conçu et déployé une architecture composée de six éléments majeurs : (1)~une \textbf{infrastructure à clés publiques (PKI)} basée sur HashiCorp Vault avec une hiérarchie CA racine -- CA intermédiaire ; (2)~un \textbf{broker MQTT Mosquitto} configuré en TLS~1.3 et mTLS obligatoire ; (3)~un \textbf{simulateur de flotte IoT} rejouant 76\,000 événements réels issus d'un dataset CSV couvrant sept types d'attaques ; (4)~une \textbf{analyse réseau} par Suricata avec règles personnalisées ; (5)~un \textbf{SIEM Wazuh} ingérant les alertes et corrélant les événements ; (6)~un \textbf{pipeline de détection d'anomalies} par Machine Learning (IsolationForest + RandomForest).

Les résultats obtenus démontrent l'efficacité de l'approche : le RandomForest atteint un F1-score de \textbf{0,994} et un ROC-AUC de \textbf{0,999} pour la détection binaire, tandis que l'IsolationForest, en mode non supervisé, obtient \textbf{0,930} / \textbf{0,959}. Les tests de performance montrent que l'overhead TLS~1.3/mTLS sur le RTT MQTT est de \textbf{+89\,\%} (4,2\,ms $\to$ 7,9\,ms), restant largement en deçà du seuil critique de 50\,ms pour les applications IoT industrielles.

\vspace{0.5cm}
\hrule
\vspace{0.3cm}

\textbf{Abstract} -- \textit{The growth of IoT in telecom operator environments generates massive and heterogeneous data flows exposed to increasing risks of interception, alteration and compromise. This final-year project, carried out in partnership with Tunisie Telecom and following Scrum methodology, addresses the end-to-end security of communication flows in a simulated IoT ecosystem. We designed and deployed a complete architecture including a PKI based on HashiCorp Vault, a Mosquitto MQTT broker with TLS~1.3/mTLS, an IoT fleet simulator, Suricata IDS, Wazuh SIEM and an ML anomaly detection pipeline. The RandomForest classifier achieves F1=0.994, ROC-AUC=0.999. TLS overhead on RTT is +89\,\% (4.2\,ms to 7.9\,ms), well within IoT industrial thresholds.}

\cleardoublepage
